% Tutorial set: 02
% Source: T2-CSP Game Playing, p. 1, Questions 2.1--2.2
% Reference: T2-CSP Game Playing-solution, pp. 1--6
% Session date: 2026-09-11
% Human verified: Search counts, heuristic values, minimax backups, and caveats checked

\subsection{Tutorial 02：CSP 与 Game Playing}
\label{tutorial:SC3000-TUT02}

\exercisemeta
  {SC3000-TUT02}
  {2026-09-11}
  {T2-CSP Game Playing，p.~1，Questions 2.1--2.2；官方参考解答 pp.~1--6}
  {搜索规模、heuristic values 与 Minimax backup 均已核对；参考答案的公平性结论已订正}
  {已确认（2026-09-11）}

\exercisequestion{SC3000-TUT02-Q01}{Logo Coloring：Forward Checking 与 CSP Heuristics}

\textbf{对应讲义：}
\begin{itemize}
  \item \relatedtopic{SC3000-L02-T04}{CSP Formulation 与 Backtracking}
  \item \relatedtopic{SC3000-L02-T05}{CSP Heuristics 与 Min-Conflicts}
\end{itemize}

\subsubsection*{题目}

Intelligent Systems Laboratory 的 logo 由 $D,I,S,L$ 四个 design elements 构成。每个 element 可染成 red、green 或 blue，且相邻 elements 必须颜色不同。使用带 Forward Checking（前向检查）的 Depth-First Search（深度优先搜索）生成所有合法 color combinations。分别完成：

\begin{enumerate}[label=(\alph*)]
  \item 按 alphabetical order（字母顺序）$D\to I\to L\to S$ 赋值；
  \item 使用适合 CSP 的最佳 variable/value heuristics 决定赋值顺序，并说明其他 heuristics 是否有用。
\end{enumerate}

两种情形均需给出完整 search-space profile（搜索空间各层规模）、total nodes、average branching factor，并评价效率。

\begin{checkedsolution}
令 variables、domains 与 constraints 为
\[
  X=\{D,I,S,L\},\qquad
  \operatorname{Dom}(X_i)=\{R,G,B\},
\]
\[
  D\ne I,\qquad D\ne S,\qquad I\ne S,\qquad S\ne L.
\]
约束图是 triangle $D-I-S$ 加上一条 tail $S-L$。Forward checking 在赋值后立即从尚未赋值的邻居 domains 中删除该颜色；若某个 domain 变成空集，则立即 backtrack。

\begin{figure}[htbp]
  \centering
  \resizebox{0.98\textwidth}{!}{\begin{tikzpicture}[
  >=Latex,
  graphnode/.style={circle,draw=noteBlue,fill=noteBlue!8,thick,minimum size=8mm,font=\small\bfseries},
  levelnode/.style={rounded corners,draw=gray!70,fill=gray!7,minimum width=26mm,minimum height=7mm,align=center,font=\scriptsize},
  every path/.style={thick}
]
  \node[font=\small\bfseries] at (-5.6,2.65) {Constraint graph};
  \node[graphnode] (D) at (-6.6,1.35) {$D$};
  \node[graphnode] (I) at (-5.2,1.35) {$I$};
  \node[graphnode] (S) at (-5.9,0.0) {$S$};
  \node[graphnode] (L) at (-4.4,0.0) {$L$};
  \draw (D)--(I) (D)--(S) (I)--(S) (S)--(L);
  \node[font=\scriptsize,align=center] at (-5.55,-1.05)
    {$\deg(S)=3$，$\deg(D)=\deg(I)=2$，$\deg(L)=1$};

  \node[font=\small\bfseries] at (-0.9,2.65) {Alphabetical: $D\to I\to L\to S$};
  \node[levelnode] (a0) at (-0.9,1.95) {root: $1$};
  \node[levelnode] (a1) at (-0.9,1.15) {$D$: $3$};
  \node[levelnode] (a2) at (-0.9,0.35) {$I$: $6$};
  \node[levelnode] (a3) at (-0.9,-0.45) {$L$: $18$};
  \node[levelnode] (a4) at (-0.9,-1.25) {$S$: $12$ solutions};
  \draw[->] (a0)--(a1); \draw[->] (a1)--(a2); \draw[->] (a2)--(a3); \draw[->] (a3)--(a4);

  \node[font=\small\bfseries] at (4.05,2.65) {Degree: $S\to D\to I\to L$};
  \node[levelnode] (h0) at (4.05,1.95) {root: $1$};
  \node[levelnode] (h1) at (4.05,1.15) {$S$: $3$};
  \node[levelnode] (h2) at (4.05,0.35) {$D$: $6$};
  \node[levelnode] (h3) at (4.05,-0.45) {$I$: $6$};
  \node[levelnode] (h4) at (4.05,-1.25) {$L$: $12$ solutions};
  \draw[->] (h0)--(h1); \draw[->] (h1)--(h2); \draw[->] (h2)--(h3); \draw[->] (h3)--(h4);
\end{tikzpicture}
}
  \caption{Logo CSP 的 constraint graph 与两种顺序的 search-space profile（原创重绘）。}
  \label{fig:sc3000-tut02-csp-profile}
\end{figure}

\paragraph{(a) Alphabetical order.}
各层 node counts（包括 root）为
\[
  \boxed{1,\ 3,\ 6,\ 18,\ 12},
  \qquad N=\boxed{40}.
\]
$D$ 与 $I$ 必须使用不同颜色。$L$ 暂时有三种选择，但其中一种会用掉 $D,I$ 之外的第三种颜色，使 $S$ 的 domain 变空；所以 18 个 $L$ nodes 中只有 12 个能生成 $S$ solution node。

以实际产生 children 的 nodes 为 internal nodes，tree edges 为 $N-1$，故
\[
  b_{\mathrm{avg}}
  =\frac{N-1}{N_{\mathrm{internal}}}
  =\frac{39}{1+3+6+12}
  =\boxed{\frac{39}{22}\approx1.77}.
\]
相较不使用 forward checking 的完整三叉树 $1+3+3^2+3^3+3^4=121$ nodes，搜索量已明显下降，但把约束最多的 $S$ 留到最后仍产生无效的 $L$ assignments。

\paragraph{(b) Heuristic ordering.}
Degree heuristic（度启发式）选择与最多 unassigned variables 相邻的 variable：
\[
  \deg(S)=3,\quad \deg(D)=\deg(I)=2,\quad \deg(L)=1.
\]
因此可取 $S\to D\to I\to L$（交换 $D,I$ 等价）。各层为
\[
  \boxed{1,\ 3,\ 6,\ 6,\ 12},
  \qquad N=\boxed{28},
\]
\[
  b_{\mathrm{avg}}
  =\frac{27}{1+3+6+6}
  =\boxed{\frac{27}{16}\approx1.69}.
\]
MRV / Most Constrained Variable（最少剩余值）在初始时因四个 domains 都有三个值而打平；LCV / Least Constraining Value（最少约束值）也因三种颜色完全对称而不能区分首选颜色。Degree heuristic 能直接选出 $S$，将 40 nodes 进一步降至 28，同时保留全部 $3!\times2=12$ 个合法 colorings。
\end{checkedsolution}

\textbf{检查点：}$N-1$ 是 tree edges 数，不是 total nodes；average branching factor 应除以实际 expanded internal nodes，而不是 tree depth。

\exercisequestion{SC3000-TUT02-Q02}{Tic-Tac-Toe：Symmetry、Evaluation 与 Minimax Backup}

\textbf{对应讲义：}
\begin{itemize}
  \item \relatedtopic{SC3000-L03-T01}{Game Formulation 与 Minimax}
  \item \relatedtopic{SC3000-L03-T02}{Cutoff、Evaluation、Quiescent Search 与 MCTS}
\end{itemize}

\subsubsection*{题目}

Tic-Tac-Toe（井字棋）由 X 与 O 轮流在 $3\times3$ board 的空格落子；同一 row、column 或 diagonal 上先连成三个相同 marks 的玩家获胜。定义 $X_n$ 为“恰有 $n$ 个 X 且没有 O”的 winning lines 数量，$O_n$ 类似。对 non-terminal position 使用
\[
  h(s)=3X_2+X_1-3O_2-O_1;
\]
X 胜、O 胜与 draw 的 terminal utilities 分别为 $+1,-1,0$。

\begin{enumerate}[label=(\alph*)]
  \item 利用 symmetry（对称性）画出 depth $2$ 的 game tree，即 X、O 各走一步；
  \item 计算所有 depth-2 positions 的 heuristic values，使用 Minimax 回传 depth $1$ 与 root values，并确定 X 的最佳第一步。
\end{enumerate}

\begin{checkedsolution}
空 board 的九个位置在 rotations/reflections 下分为 corner、centre、edge 三类，所以 depth $1$ 只需保留三个不等价 moves。O 的回应分别有 $5,2,5$ 个 symmetry classes，depth $2$ 共 $12$ 个不等价 positions。

在 depth $2$ 只有一个 X 和一个 O，因此
\[
  X_2=O_2=0,
  \qquad
  \boxed{h(s)=X_1-O_1}.
\]
同时含 X 与 O 的 line 已被 block，不计入任何一方的 $X_1$ 或 $O_1$。

\begin{figure}[htbp]
  \centering
  \resizebox{0.98\textwidth}{!}{\begin{tikzpicture}[
  >=Latex,
  edge/.style={draw=gray!75,-{Latex[length=2mm]},thick},
  maxnode/.style={circle,draw=noteBlue,fill=noteBlue!10,thick,minimum size=10mm,align=center,font=\scriptsize},
  minnode/.style={rounded corners,draw=ntuRed,fill=ntuRed!8,thick,minimum width=29mm,minimum height=11mm,align=center,font=\scriptsize},
  leaf/.style={circle,draw=gray!75,fill=gray!7,minimum size=7mm,font=\scriptsize}
]
  \node[maxnode] (root) at (0,3.6) {MAX\\$+1$};
  \node[minnode] (corner) at (-5.0,2.15) {corner\\MIN $=-1$};
  \node[minnode] (centre) at (0,2.15) {centre\\MIN $=+1$};
  \node[minnode] (edge) at (5.0,2.15) {edge\\MIN $=-2$};
  \draw[edge] (root)--(corner); \draw[edge] (root)--(centre); \draw[edge] (root)--(edge);

  \foreach \x/\v [count=\i] in {-7/1,-6/-1,-5/0,-4/0,-3/1}{
    \node[leaf] (c\i) at (\x,0.35) {$\v$};
    \draw[edge] (corner)--(c\i);
  }
  \foreach \x/\v [count=\i] in {-0.55/1,0.55/2}{
    \node[leaf] (m\i) at (\x,0.35) {$\v$};
    \draw[edge] (centre)--(m\i);
  }
  \foreach \x/\v [count=\i] in {3/-1,4/-2,5/0,6/-1,7/0}{
    \node[leaf] (e\i) at (\x,0.35) {$\v$};
    \draw[edge] (edge)--(e\i);
  }

  \node[font=\scriptsize,align=center] at (0,-0.55)
    {Leaf values use $h=3X_2+X_1-3O_2-O_1=X_1-O_1$ at depth $2$};
\end{tikzpicture}
}
  \caption{利用 symmetry 压缩后的 depth-2 game tree。Leaf values 来自 $h=X_1-O_1$；MIN 与 MAX 的回传值标在内部节点中（原创重绘）。}
  \label{fig:sc3000-tut02-tictactoe-tree}
\end{figure}

三类第一步的 leaf values 与 backup 为
\begin{center}
\small
\begin{tabularx}{0.94\textwidth}{@{}lXcc@{}}
  \toprule
  X first move & Depth-2 heuristic values & MIN backup & Root comparison \\
  \midrule
  corner & $1,-1,0,0,1$ & $-1$ & \\
  centre & $1,2$ & $+1$ & $\leftarrow$ maximum \\
  edge & $-1,-2,0,-1,0$ & $-2$ & \\
  \bottomrule
\end{tabularx}
\end{center}
因此 root 的 MAX backup 为
\[
  \max\{-1,+1,-2\}=\boxed{+1},
\]
在本题的 depth-2 cutoff 与 evaluation function 下，最佳第一步为
\[
  \boxed{\text{play in the centre}}.
\]

该结论只表示 centre 在当前 shallow evaluation（浅层评估）下的 worst-case value 最高，不能证明 X 一定获胜。标准 Tic-Tac-Toe 在双方 optimal play 下是 draw；参考答案所称 ``whoever plays first has a strong advantage'' 不能作为 game-theoretic forced win 的结论。

此外，题设把 terminal win 记为 $+1$，但 non-terminal heuristic 可超过 $+1$。更稳妥的实现应令 terminal win/loss 使用超过所有 heuristic values 的 $\pm M$（或 $\pm\infty$），避免算法把未结束局面的 $+2$ 错排在真实胜利 $+1$ 之前。
\end{checkedsolution}

\textbf{检查点：}Minimax 的 MAX/MIN backup 给出的是当前 search depth 与 evaluation function 下的决策；只有搜索到 terminal states，或 evaluation 足够准确时，才可把回传值解释为最终胜负保证。

\subsubsection*{本次 Tutorial 收口}

\begin{itemize}
  \item Q2.1：Forward checking 通过删减未来 domains 剪枝；Degree heuristic 优先处理 constraint graph 的高连接变量；
  \item Q2.1：区分 total nodes、tree edges、internal nodes 与 average branching factor；
  \item Q2.2：先用 symmetry 合并等价 moves，再在 cutoff leaves 计算 evaluation，MIN 取最小、MAX 取最大；
  \item Q2.2：``当前 heuristic 下最优'' 不等于 ``一定获胜''，terminal utility 的尺度还应压过所有 non-terminal estimates。
\end{itemize}
